\documentclass[11pt,a4paper]{article}
\usepackage[margin=24mm]{geometry}
\usepackage{amsmath,amssymb,amsthm,hyperref}
\setlength{\emergencystretch}{1em}
\newtheorem{theorem}{Theorem}
\newtheorem{lemma}[theorem]{Lemma}
\newtheorem{proposition}[theorem]{Proposition}
\newcommand{\F}{\mathbb F}
\newcommand{\ord}{\operatorname{ord}}
\title{Square patterns and cyclotomic degree bounds over finite fields}
\author{Human Compression project: F011 common theory}
\date{}
\begin{document}\maketitle
This paper reworks the character-sum and cyclotomic mechanism of
Kim--Yip--Yoo~\cite{KYY}, with explicit zero and order exclusions.
Throughout, $q$ is a power of an odd prime $p$, the quadratic character
$\chi$ is extended by $\chi(0)=0$, and ``square'' includes zero.
We retain the classical multiplicative Weil estimate: if a nontrivial
character has order $d$, and a polynomial is not a scalar multiple of a
$d$th power and has $h$ distinct geometric roots, its complete character
sum has absolute value at most $(h-1)\sqrt q$.
This is source Lemma~2.1, citing Lidl--Niederreiter, Theorem~5.41. It is an
external theorem, not proved by this paper. The project's earlier
character-sum audit used the same interface for a different
application and did not recover the Diophantine-tuple results.

\section{Character indicators, including their zeros}
\begin{lemma}[Prescribed character values]\label{patterns}
Let $d\mid q-1$, let $\psi$ have order $d>1$, and let $a_1,\ldots,a_k$
be distinct. If $\epsilon_i$ are $d$th roots of unity and $N$ counts the
solutions of $\psi(x+a_i)=\epsilon_i$ for all $i$, then
\[
 \left|N-\frac q{d^k}\right|
 \le \left(k-1-\frac kd+\frac1{d^k}\right)\sqrt q+\frac kd.
\]
\end{lemma}
\begin{proof}
Use
$P_i(x)=d^{-1}\sum_{j=0}^{d-1}\epsilon_i^{-j}\psi(x+a_i)^j$,
where the zeroth term is one even at a root. Off $-a_i$, this is the
indicator of the desired character value; at $-a_i$ it is $1/d$.
Hence the sum of $\prod P_i$ differs from $N$ by a number between zero
and $k/d$. A nonconstant term in the expansion has exponents between
zero and $d-1$, so its linear-factor product is not a $d$th power.
For a support of size $h$ the Weil bound is $(h-1)\sqrt q$.
The total coefficient is
\[
 d^{-k}\sum_{h=1}^k\binom kh(d-1)^h(h-1)
 = k-1-k/d+d^{-k},
\]
by differentiating $(1+z)^k$ and subtracting its nonconstant part.
The triangle inequality proves the result.
\end{proof}

\begin{lemma}[Weighted square filter]\label{weighted}
Let $f_1,\ldots,f_T$ be nonconstant, monic, squarefree, pairwise coprime
polynomials, of total degree $D$. Set
\[
 W(x)=2^{-T}\prod_{i=1}^T(1+\chi(f_i(x))).
\]
Then $0\le W(x)\le1$, and $W(x)>0$ precisely when all $f_i(x)$ are
squares, possibly zero. Moreover
\[
 \sum_{x\in\F_q}W(x)\ge
 \frac q{2^T}-\left(\frac D2-1+2^{-T}\right)\sqrt q.
 \tag{1}
\]
\end{lemma}
\begin{proof}
Every factor is zero, one or two. Every nonempty subproduct of the $f_i$
is squarefree of positive degree, so the Weil bound applies. Each $f_i$
occurs in $2^{T-1}$ nonempty subproducts. Their total degree is therefore
$2^{T-1}D$, and subtracting one for each of the $2^T-1$ subproducts gives
the coefficient in (1). The filter is generally a weighted count at zeros,
not the cardinality of the admissible set; positivity suffices below.
\end{proof}

\section{Cyclotomic factors and a degree estimate}
Let $\Phi_j$ denote the reduction of the integral cyclotomic polynomial.
For $j=p^a b$, $p\nmid b$, $a\ge1$, the integral cyclotomic identities
and Frobenius give
\[
 \Phi_j(X)=\Phi_b(X)^{(p-1)p^{a-1}}\quad\hbox{in }\F_q[X]. \tag{2}
\]
This is a square. If $p\nmid j$, the polynomial $X^j-1$ is separable,
so $\Phi_j$ is squarefree. Distinct prime-to-$p$ indices give coprime
factors of the separable polynomial $X^{\operatorname{lcm}(i,j)}-1$.
None of these cyclotomic polynomials vanishes at zero. These statements
recover all three clauses of source Lemma~2.4.

For $m\ge2$, put
\[
 I=\{1\le i<m:p\nmid i\},\quad T=|I|+1=m-\lfloor(m-1)/p\rfloor,
 \quad \alpha=\frac p{p-1},
\]
and $s=\sum_{i\in I}\varphi(i)$, $S=\sum_{i\in I}\varphi(2i)$.
\begin{lemma}\label{degree}
The two degree sums satisfy
\[
 s\le S\le \sum_{i\in I}i\le K:=\frac\alpha2 T(T-1).
 \tag{3}
\]
In particular $S<\alpha T^2$, as asserted in source Lemma~3.4.
\end{lemma}
\begin{proof}
For odd $i$, $\varphi(2i)=\varphi(i)\le i$; for even $i$,
$\varphi(2i)=2\varphi(i)\le i$. Write $m-1=pu+v$, $0\le v<p$.
Then $r=|I|=(p-1)u+v=T-1$, and a sum over complete blocks gives
\[
 A:=\sum_{i\in I}i=\frac{p(p-1)u^2}2+puv+\frac{v(v+1)}2.
\]
Direct subtraction gives
\[
 \frac\alpha2r(r+1)-A=\frac{pu}2+\frac{v(v+1)}{2(p-1)}\ge0.
\]
This proves (3). It does not require Cohen's totient-sum estimate.
\end{proof}

\begin{proposition}[Two cyclotomic square patterns]\label{existence}
Suppose $m\ge2$ and
\[
 \sqrt q>\alpha\,2^{T-2}T^2. \tag{4}
\]
There is a square $y\in\F_q^*$ of order at least $m$ for which every
$y^i+1$, $1\le i<m$, is a square. If $q\equiv1\pmod4$, there is also
such a $y$ for which every $y^i-1$, $1\le i<m$, is a square.
These are separate choices of $y$.
\end{proposition}
\begin{proof}
For the plus pattern use the pairwise coprime squarefree family
$X,\Phi_{2i}(X)$ for $i\in I$. Its total degree is $D=1+S\le1+K$;
the degree-one factor must be included. Apply Lemma~\ref{weighted}.
There are at most $s$ nonzero elements of order below $m$, since for
each possible order $j$ there are at most $\varphi(j)$ elements, and
$p\nmid j$. Each removes weight at most one. The point zero has weight
$1/2$, since $\Phi_{2i}(0)=1$. Thus the weight at eligible points is at least
\[
 \frac q{2^T}-\left(\frac{K-1}2+2^{-T}\right)\sqrt q-K-\frac12.
 \tag{5}
\]
Since $T\ge2$, we have $K+1\le\alpha T^2/2$. Condition (4) gives
$z:=\sqrt q>2^{T-1}(K+1)$. Substituting this into the first term in
(5), the expression is greater than
\[
 (1-2^{-T})z-K-\tfrac12
 > (2^{T-1}-\tfrac32)K+2^{T-1}-1>0.
\]
Choose an eligible point of positive weight. It is nonzero and square;
all $\Phi_{2i}(y)$ with $i\in I$ are squares. The polynomial identity
\[
 X^j+1=\prod_{d\mid2j,\ d\nmid j}\Phi_d(X)
\]
shows the same for $y^j+1$: all such $d$ are even and at most $2j$,
and factors with $p\mid d$ are squares by (2). This is a polynomial
identity and needs no division by a possibly vanishing $y^j-1$.

For the minus pattern use $X,\Phi_i(X)$ for $i\in I$. Its degree is
$1+s\le1+K$. At zero its weight is again $1/2$ because the only negative
constant term is $\Phi_1(0)=-1$, a square under the hypothesis.
The same argument applies, using $X^j-1=\prod_{d\mid j}\Phi_d(X)$.
\end{proof}

\begin{lemma}[Explicit parameter choice]\label{parameter}
Let $q>7$, put
\[
 B=\frac{\frac12\log q-2\log\log q}{\log2},\qquad
 m=\left\lfloor\alpha(B+1)\right\rfloor.
\]
If $m\ge2$, condition (4) holds.
\end{lemma}
\begin{proof}
Writing $m=pu+v$, $0\le v<p$, directly gives $T\le m/\alpha+1$,
including the case $v=0$. Consequently $T\le B+2$.
Set $x=\log q>2$. Since $\log x>\log2$,
\[
 T\le\frac{x/2-2\log x}{\log2}+2<\frac{x}{2\log2}.
\]
Also $\alpha\le3/2<4(\log2)^2$. One elementary lower bound
$\log2=\int_1^2t^{-1}dt>2/3$ follows from strict convexity and the
midpoint rule; hence $4(\log2)^2>16/9>3/2$.
Taking logarithms, we obtain
\[
 \log\bigl(\alpha2^{T-2}T^2\bigr)
 \le x/2-2\log x+\log\alpha+2\log T<x/2,
\]
which is (4).
\end{proof}

\section{A shorter pattern with a stronger order exclusion}
\begin{proposition}\label{centered-pattern}
For an integer $n\ge4$, define
\[
 I_n=\{1\le i\le n-2:p\nmid i\},\qquad
 T=n-1-\lfloor(n-2)/p\rfloor.
\]
If $\sqrt q>\alpha2^{T-2}T^2$, there is a square $y\in\F_q^*$
of order at least $n$ such that $y^j+1$ is a square for
$1\le j\le n-2$.
\end{proposition}
\begin{proof}
Apply Lemma~\ref{degree} with $m=n-1$, and set
$K=\alpha T(T-1)/2$. The family $X,\Phi_{2i}(X)$ for $i\in I_n$
has degree at most $1+K$. This time we exclude all orders below $n$,
not merely those below $n-1$. Their number is at most
$s+\varphi(n-1)\le K+n-1\le K+\alpha T$; the last inequality follows
from $T\ge(n-1)/\alpha$. Exclude zero at weight $1/2$ as before.
Writing $z=\sqrt q$, the remaining weight is at least
\[
 \frac {z^2}{2^T}
 -\left(\frac{K-1}{2}+2^{-T}\right)z-K-\alpha T-\frac12.
 \tag{6}
\]
Here $T\ge3$, so $K\ge\alpha T$. Also
$K+1\le\alpha T^2/2$, and the hypothesis gives
$z>2^{T-1}(K+1)$. Thus (6) is greater than
\[
 (2^{T-1}-\tfrac32)K-\alpha T+2^{T-1}-1
 \ge(2^{T-1}-\tfrac52)K+2^{T-1}-1>0.
\]
The same cyclotomic product identity as in
Proposition~\ref{existence} now gives every required power condition.
The extra order exclusion is essential for the centered construction.
\end{proof}

\begin{lemma}[Parameter choice for the stronger bound]\label{parameter-plus-two}
Suppose $q>7$ and $q\equiv7\pmod8$. Put
\[
 B=\frac{\frac12\log q-2\log\log q}{\log2},\qquad
 n=\lfloor\alpha(B+2)\rfloor.
\]
If $n\ge4$, the hypothesis of Proposition~\ref{centered-pattern} holds.
\end{lemma}
\begin{proof}
The characteristic satisfies $p\equiv7\pmod8$, the extension degree is
odd, and $p\ge7$. Write $n=pu+v$, $0\le v<p$.
If $v=0$, then $T=(p-1)u=n/\alpha$; if $v=1$, then
$T=(p-1)u+1=n/\alpha+1/p$; and if $v\ge2$, then
$T=(p-1)u+v-1=n/\alpha+v/p-1$. Therefore
\[
 T\le n/\alpha+1/p\le B+2+1/p. \tag{7}
\]
This replaces the insufficient rounding estimate in the source proof.

Set $x=\log q$. The smallest eligible prime power beyond seven is
at least $23$, so $x>3$. Since $2+1/p\le15/7<3$ and $9>8$,
we have $2\log x>(2+1/p)\log2$. Consequently (7) implies
$T<x/(2\log2)$. Moreover
\[
 \alpha2^{1/p}\le\tfrac76\,2^{1/7}
 <\tfrac74<\tfrac{16}{9}<4(\log2)^2;
\]
here $2^{1/7}<3/2$ and the final inequality was proved in
Lemma~\ref{parameter}. Hence
\begin{align*}
 \log(\alpha2^{T-2}T^2)
 &\le x/2-2\log x+\log(\alpha2^{1/p})+2\log T\\
 &<x/2.
\end{align*}
This proves the stated strict inequality.
\end{proof}

\begin{thebibliography}{9}
\bibitem{KYY} S. Kim, C. H. Yip and S. Yoo,
\emph{Explicit constructions of Diophantine tuples over finite fields},
Ramanujan J. \textbf{65} (2024), 163--172,
doi:10.1007/s11139-024-00888-5.
\bibitem{LN} R. Lidl and H. Niederreiter, \emph{Finite Fields}, second
edition, Cambridge University Press, 1997, Theorem~5.41 (retained through
the exact statement in \cite[Lemma 2.1]{KYY}).
\end{thebibliography}
\end{document}
